% ---------------------------------------------------------------------
%  Chapter 10 : From Bare Graphs to Labelled Diagrams
%               (Type assignment, causality, symmetry factor, filter)
% ---------------------------------------------------------------------
\chapter{From Bare Graphs to Labelled Diagrams}\label{ch:type-assignment}

The previous chapter handed us a pile of \term{prediagrams}: bare directed
multigraphs with $k$ external leaves and $\ell$ loops, deduplicated up to graph
isomorphism, but with \emph{no physics on them yet}. A prediagram knows that
``two arrows leave this vertex and one enters it''; it does not yet know which
monomial of the action sits there, which field rides each leg, or which
propagator each edge carries. This chapter is the bridge from those bare
topologies to fully labelled Feynman diagrams that an integrator can actually
evaluate.

The bridge is four pipeline phases, run in order, and they are the subject of
this chapter:

\begin{description}
  \item[Phase D --- Filter] (\file{msrjd/diagrams/filter.py}). A cheap
        set-membership pre-screen that throws away any prediagram whose vertex
        \emph{degrees} no vertex or source of the theory could ever supply ---
        before we pay for the expensive labelling.
  \item[Phase E --- Type assignment] (\file{msrjd/diagrams/type\_assignment.py}).
        The heart of the chapter. Given one bare graph, produce \emph{every}
        valid way to decorate it: a monomial at each internal vertex, an external
        field on each leaf, and a propagator $G_{ij}$ on each edge. This is a
        constraint-satisfaction search; its output is a stream of
        \code{TypedDiagram} objects.
  \item[Phase F --- Causality] (\file{msrjd/diagrams/causality.py}). Discard typed
        diagrams that violate the retarded boundary condition of \msrjd{}.
        Structurally, that means the diagram must be a directed acyclic graph.
  \item[Phase G --- Symmetry factor \& dedup]
        (\file{msrjd/diagrams/symmetry.py}). Deduplicate the typed diagrams (Phase
        E emits the same diagram many times, relabelled) and compute the
        combinatorial \emph{symmetry factor} $\Sym(\Gamma)$ that every Feynman
        rule carries.
\end{description}

A fifth file, \file{msrjd/fork\_safety.py}, is a small shared safety guard that
the parallel branch of Phase E consults; we treat it at the end of the chapter
because it is shared with the temporal Phase~J path and is one of the most
load-bearing six lines in the whole repository.

\begin{note}
A quick orientation on \emph{why these phases are separate}. Filtering is
$O(\text{vertices})$ and uses no algebra; type assignment is a backtracking
search; deduplication and the symmetry factor both need a graph-isomorphism
engine (\code{nauty}, via Sage). Keeping them apart lets the cheap check run
first, lets the search be tested in isolation, and confines every call into the
isomorphism machinery to one file. The single function that wires all four
together with disk caching is
\code{enumerate\_unique\_diagrams} in \file{pipeline/\_diagrams.py}; we read it in
\S\ref{sec:ta-orchestrator}.
\end{note}

% =====================================================================
\section{The two field species and the edge convention}\label{sec:ta-convention}

Everything in this chapter rests on one convention, so we state it before any
code. A reader fluent in \msrjd{} field theory will recognise all of it; we
restate it from the ground up so the code's index choices are unambiguous.

\subsection{Physical fields and response fields}

A Langevin equation $\Dt\varphi = F[\varphi] + \xi$ with noise $\xi$ is rewritten,
\`a la Martin--Siggia--Rose / Janssen--De~Dominicis (\msrjd{}), as a path integral
over \emph{two} fields per physical field. The generating functional is
\begin{equation}
  Z \;=\; \int \mathcal{D}[\varphi]\,\mathcal{D}[\tilde\varphi]\;
  \ee^{-S[\varphi,\,\tilde\varphi]} .
  \label{eq:ta-Z}
\end{equation}
The two species are:

\begin{itemize}
  \item the \term{physical field} $\varphi$ --- the actual dynamical variable. In
        the code's naming its generator base names are letters like \code{vt},
        \code{nt}, \code{xt} (the trailing ``t'' is part of the base name, not a
        tilde);
  \item the \term{response field} $\tilde\varphi$ (``tilde'') --- the auxiliary
        conjugate field of \msrjd{}. In the code its base names carry a \code{d}
        prefix: \code{dv}, \code{dn}, \code{dx}.
\end{itemize}

The quadratic (Gaussian) part of $S$ defines the \emph{propagators}; the
higher-order part defines the \emph{interaction vertices}; and the noise enters
as \emph{source vertices} --- pure response-field monomials drawn from the noise
cumulant-generating functional.

\subsection{What a directed edge means}

Every internal line of a diagram is a propagator, and \msrjd{} has essentially
one forward-in-time propagator family: the retarded response $\avg{\varphi\,
\tilde\varphi}$. The code encodes this with a strict directionality convention,
stated at the top of \file{type\_assignment.py:11--13}:

\begin{quote}\itshape
each directed edge $u\to v$ carries a propagator $G_{ij}$ where $i$ is the
response-field index contributed by vertex $u$ (the tail) and $j$ is the
physical-field index contributed by vertex $v$ (the head).
\end{quote}

So a directed edge is a \emph{Wick contraction} between a response leg
$\tilde\varphi_i$ at its tail and a physical leg $\varphi_j$ at its head. The
value carried is an entry of the propagator matrix
$G_{\mathrm{ft}} = K_{\mathrm{ft}}^{-1}$, the inverse of the quadratic-form matrix
$K$. Now the part that is genuinely easy to get backwards, and which the code
warns about in two places (\file{type\_assignment.py:293--296} and
\file{:490--502}):

\begin{defn}
\textbf{The $(pi, ri)$ index order.}
$G_{\mathrm{ft}} = K_{\mathrm{ft}}^{-1}$ has \emph{rows indexed by physical fields
and columns by response fields}, so
\[
  G_{\mathrm{ft}}[i,\,j] \;=\; \avg{\varphi_i\,\tilde\varphi_j}.
\]
The propagator on an edge whose \emph{tail} contributes response index $ri$ and
whose \emph{head} contributes physical index $pi$ is therefore
\[
  \avg{\varphi_{\mathrm{phys}}\,\tilde\varphi_{\mathrm{resp}}}
  \;=\; G_{\mathrm{ft}}[\,pi,\; ri\,].
\]
You index $G_{\mathrm{ft}}$ with \textbf{(physical-row, response-col)}, in that
order.
\end{defn}

This ordering is load-bearing. The code stores the pair the ``natural'' way
round --- \code{propagator\_indices[edge] = (ri, pi)} --- but every \emph{lookup}
into the matrix or its zero-mask uses \code{(pi, ri)}. We will see exactly that
asymmetry in the backtracker, and \S\ref{sec:ta-gotchas} records the historical
bug that motivated the comment.

\begin{figure}[ht]
\centering
\begin{tikzpicture}[
    >={Stealth[length=2.6mm]},
    vert/.style={circle,draw,fill=notebg,minimum size=7mm,inner sep=0pt},
    every node/.style={font=\small},
  ]
  \node[vert] (u) at (0,0) {$u$};
  \node[vert] (v) at (4.4,0) {$v$};
  \draw[->,thick] (u) -- (v)
    node[midway,above=3pt] {$G_{\mathrm{ft}}[\,pi,\,ri\,]
        = \avg{\varphi_j\,\tilde\varphi_i}$};
  \node[below=2pt,align=center] at (u.south)
    {tail: response\\leg $\tilde\varphi_i$ ($ri$)};
  \node[below=2pt,align=center] at (v.south)
    {head: physical\\leg $\varphi_j$ ($pi$)};
\end{tikzpicture}
\caption{The edge convention. A directed edge $u\to v$ is a Wick contraction of a
response leg at the tail with a physical leg at the head; its value is
$G_{\mathrm{ft}}$ indexed (physical-row, response-col).}
\label{fig:ta-edge}
\end{figure}

\subsection{Retarded causality, stated two ways}

The retarded Green's function $\Gret(t)$ is nonzero only for $t>0$. In Fourier
space that is the statement that \emph{all poles of $\det\hat K(\omega)$ lie in
the upper half of the complex $\omega$ plane}, $\operatorname{Im}\omega_k > 0$, so
that closing the contour the ``correct'' way reproduces the retarded boundary
condition. There are two equivalent manifestations, and Phase F uses the first:

\begin{itemize}
  \item \textbf{Structural.} A diagram built only from retarded responses and
        noise vertices is a \term{directed acyclic graph} (DAG): there is no way
        to return to a vertex by following time-forward arrows. A directed cycle
        would close a loop of strictly-retarded propagators, which integrates to
        zero (or signals an acausal construction). This is the operative check,
        \code{D.is\_directed\_acyclic()}.
  \item \textbf{Analytic.} If explicit pole locations are supplied, every pole
        must have $\operatorname{Im} > 0$. This lives in
        \code{check\_pole\_structure}, but in the current pipeline the typed-diagram
        path passes \emph{no} poles, so only the DAG check runs.
\end{itemize}

% =====================================================================
\section{Phase D --- the degree filter}\label{sec:ta-filter}

We begin with the cheapest phase. Phase D
(\file{msrjd/diagrams/filter.py}) never looks at a single field type; it asks
only ``could the \emph{degrees} of this graph's vertices conceivably be supplied
by the theory's monomials?'' and throws away the hopeless graphs in
$O(\text{vertices})$ time.

\subsection{Classifying a prediagram's vertices}

The first helper sorts the non-leaf vertices of a prediagram into the two kinds
of internal vertex. A \term{source} is a noise vertex: it has \emph{no} incoming
edges (in-degree $0$), because a pure response-field monomial offers only
response legs, which feed \emph{outgoing} edges. Anything else is an
\term{interaction} vertex.

\begin{lstlisting}
def classify_prediagram_vertices(D, leaves):  # filter.py:12-40
    leaf_set = set(leaves)
    source = []
    interaction = []
    for v in D.vertices():
        if v in leaf_set:
            continue
        if D.in_degree(v) == 0:
            source.append(v)
        else:
            interaction.append(v)
    return source, interaction
\end{lstlisting}

\code{D} is a Sage \code{DiGraph} (we meet Sage properly in
\S\ref{sec:ta-tools}); \code{D.in\_degree(v)} is the number of edges pointing
into $v$. The exact same source-vs-interaction split, by the same in-degree-zero
test, reappears verbatim inside Phase E
(\file{type\_assignment.py:167--177}) --- the two phases agree on what counts as a
noise vertex.

\subsection{The filter itself}

\begin{lstlisting}
def filter_prediagrams(prediagrams, vertex_types, source_types):  # filter.py:43-81
    from msrjd.core.vertices import available_degrees
    int_degs, src_degs = available_degrees(vertex_types, source_types)
    kept = []
    for pd in prediagrams:
        D, G, leaves, internal = pd
        source_verts, interaction_verts = classify_prediagram_vertices(D, leaves)
        valid = True
        for v in source_verts:
            if D.out_degree(v) not in src_degs:
                valid = False
                break
        if valid:
            for v in interaction_verts:
                if (D.in_degree(v), D.out_degree(v)) not in int_degs:
                    valid = False
                    break
        if valid:
            kept.append(pd)
    return kept, len(prediagrams) - len(kept)
\end{lstlisting}

The pivot is \code{available\_degrees(vertex\_types, source\_types)}, defined in
\file{msrjd/core/vertices.py:777}. It returns two sets:

\begin{itemize}
  \item \code{int\_degs} --- the set of $(\text{in\_degree},
        \text{out\_degree})$ \emph{pairs} that the theory's interaction vertices
        can supply, built as
        \code{\{(vt.in\_degree, vt.out\_degree) for vt in vertex\_types\}}
        (\file{vertices.py:793}). Here the \code{in\_degree} of a \code{VertexType}
        is the count of its physical legs and \code{out\_degree} the count of its
        response legs;
  \item \code{src\_degs} --- the set of \code{out\_degree} \emph{values} the
        sources can supply, \code{\{st.out\_degree for st in source\_types\}}
        (\file{vertices.py:794}).
\end{itemize}

Each prediagram passes only if \emph{every} source vertex has an out-degree in
\code{src\_degs} \emph{and} every interaction vertex has a $(\text{in},
\text{out})$ pair in \code{int\_degs}. The first failing vertex flips
\code{valid} to \code{False} and short-circuits the loop --- there is no point
checking the rest.

\begin{example}[Why this saves work]
Suppose the theory has only a cubic interaction $\varphi^2\tilde\varphi$ (so
\code{int\_degs} contains only the bidegree $(2,1)$) and a single quadratic noise
$\tilde\varphi^2$ (so \code{src\_degs}$=\{2\}$). A prediagram containing a
degree-4 internal vertex --- a quartic --- is rejected here in a handful of
integer comparisons, instead of after a fruitless backtracking search in Phase E
that would have found no candidate type for that vertex and bailed out anyway.
The filter is purely a performance gate; it never changes the final answer.
\end{example}

\begin{note}
Phase D is \emph{not} in the production orchestrator. The loop enumerator
(Phase~C, Chapter~\ref{ch:enumeration}) already emits only degree-valid
prediagrams for the chosen $(k,\ell)$, so \code{enumerate\_unique\_diagrams}
skips the explicit filter. It remains the canonical cheap pre-screen --- it is
unit-tested in \file{tests/test\_filter.py} and is the right tool for any caller
that enumerates prediagrams more loosely --- but a reader tracing the live
pipeline will not see it called.
\end{note}

% =====================================================================
\section{Phase E --- type assignment as constraint satisfaction}\label{sec:ta-typeassign}

This is the engine room. Given one bare prediagram and the theory's menu of
vertex and source types, Phase E
(\file{msrjd/diagrams/type\_assignment.py}) enumerates \emph{every} valid way to
turn the bare graph into a labelled Feynman diagram. ``Valid'' means three things
must all hold simultaneously: vertex degrees must match a real monomial, every
leg must point the right way (response legs feed outgoing edges, physical legs
feed incoming edges), and the propagator the labelling demands on each edge must
not be identically zero. It is a backtracking constraint solver, and its output
is a stream of \code{TypedDiagram} objects.

\subsection{The output object: \texttt{TypedDiagram}}

Before the algorithm, the thing it produces. A \code{TypedDiagram}
(\file{type\_assignment.py:29--71}) is a \code{\_\_slots\_\_} class --- meaning
its instances carry only the five named attributes and cannot accidentally grow
others, which keeps them memory-lean since we may make thousands. The five
fields are:

\begin{description}
  \item[\code{prediagram}] the \code{(D, G, leaves, internal)} 4-tuple this was
        built from.
  \item[\code{vertex\_assignments}] a dict
        \code{\{vertex\_id: VertexType | SourceType\}} for the internal vertices
        (empty in the no-internal case).
  \item[\code{edge\_types}] a dict
        \code{\{(u, v, label): (resp\_leg, phys\_leg)\}}. The 3-tuple edge key
        keeps parallel edges distinct (the \code{label} disambiguates them); each
        value is a pair of \code{(field\_base, pop\_idx)} legs.
  \item[\code{external\_legs}] a dict
        \code{\{leaf\_vertex: (field\_base, pop\_idx)\}}.
  \item[\code{propagator\_indices}] a dict
        \code{\{(u, v, label): (ri, pi)\}} --- the row/col into
        $G_{\mathrm{ft}}$, stored as $(ri, pi)$ (recall the lookup flips this to
        $(pi, ri)$).
\end{description}

It also defines \code{\_\_getstate\_\_}/\code{\_\_setstate\_\_}
(\file{type\_assignment.py:60--65}) so that a \code{\_\_slots\_\_} object can be
pickled across the fork pool (\S\ref{sec:ta-fork}), and a \code{\_\_repr\_\_}
that summarises vertex / edge / assignment counts.

\subsection{A field-leg is a \texttt{(base, pop\_idx)} tuple}

Throughout, a single field-leg identity is a tuple \code{(field\_base,
pop\_idx)}, e.g. \code{('nt', 1)} or \code{('dn', 2)}. The \code{field\_base} is
the generator's base name (letters); the \code{pop\_idx} is a 1-based
\term{population index} parsed off the trailing digits (the \code{2} in
\code{nt2}), distinguishing populations or components. Index \code{0} means ``no
trailing digits.'' Whether a leg is a response or physical leg is \emph{not}
encoded in the tuple --- it is decided by which generator list the base name came
from, which is exactly what the next helper records.

\subsection{The bridge from legs to matrix indices}

A leg \code{(base, pop\_idx)} must be turned into a row or column of
$G_{\mathrm{ft}}$. \code{build\_field\_index\_map}
(\file{type\_assignment.py:76--110}) builds that bridge once:

\begin{lstlisting}
def build_field_index_map(ring_var_names, n_tilde):  # type_assignment.py:76-110
    from msrjd.core.vertices import _parse_field_name
    resp_index = {}
    phys_index = {}
    resp_counter = 0
    phys_counter = 0
    for i, name in enumerate(ring_var_names):
        base, pop_idx = _parse_field_name(name)
        if i < n_tilde:
            resp_index[(base, pop_idx)] = resp_counter
            resp_counter += 1
        else:
            phys_index[(base, pop_idx)] = phys_counter
            phys_counter += 1
    return resp_index, phys_index
\end{lstlisting}

\code{ring\_var\_names} is the ordered list of polynomial-ring generator names,
e.g. \code{['vt1','vt2','nt1','nt2','dv1','dv2','dn1','dn2']}, and \code{n\_tilde}
is the number of response generators --- crucially, the response generators come
\emph{first}. So a generator at position \code{i < n\_tilde} is a response field
(it gets the next \code{resp\_counter}); otherwise it is a physical field. The
helper \code{\_parse\_field\_name} (from \file{core/vertices.py:327}) splits a
name like \code{'nt1'} into \code{('nt', 1)}. The two returned dicts,
\code{resp\_index} and \code{phys\_index}, are the maps every later phase uses to
turn a leg into a matrix index.

\subsection{The top-level enumerator}

\code{enumerate\_typed\_diagrams} (\file{type\_assignment.py:115--266}) is the
top-level generator for \emph{one} prediagram. Reading its body in order is the
clearest way to understand the algorithm; it proceeds in seven steps.

\paragraph{Step 1 --- the propagator zero mask
(\file{:158--165}).} Before searching, it precomputes which entries of
$G_{\mathrm{ft}}$ are identically zero, so the inner loop can prune any labelling
that would place a zero propagator:

\begin{lstlisting}
g_zero_mask = {}
for _i in range(G_ft.nrows()):
    for _j in range(G_ft.ncols()):
        g_zero_mask[(_i, _j)] = (str(G_ft[_i, _j]) == '0')   # type_assignment.py:165
\end{lstlisting}

The test is \code{str(G\_ft[i,j]) == '0'} --- a deliberately \emph{structural,
string-based} test, \emph{not} the semantically-correct \code{is\_zero()}. The
reason is performance, and it matters enough that \S\ref{sec:ta-zeromask} is
devoted to it. For now: this catches only entries that are the literal
\code{SR(0)}, which is exactly what the propagator builder writes into empty
cells.

\paragraph{Step 2 --- classify the non-leaf vertices
(\file{:167--177}).} The same in-degree-zero split as Phase D, producing
\code{source\_verts} and \code{interaction\_verts}; their concatenation
\code{ordered\_internal = source\_verts + interaction\_verts} fixes a
deterministic processing order for the backtracker.

\paragraph{Step 3 --- per-vertex incident-edge lists (\file{:179--184}).} For
every vertex, cache \code{D.outgoing\_edges(v)} and \code{D.incoming\_edges(v)}.
These come back as Sage 3-tuples \code{(u, v, label)} so that parallel edges stay
distinct.

\paragraph{Step 4 --- candidate types per vertex (\file{:186--201}).} A source
vertex of out-degree \code{od} gets every \code{SourceType} with
\code{st.out\_degree == od}; an interaction vertex of bidegree \code{(ind, od)}
gets every \code{VertexType} whose \code{(in\_degree, out\_degree)} matches. If
\emph{any} vertex has \emph{no} candidate, the whole prediagram is impossible and
the generator simply \code{return}s, yielding nothing.

\paragraph{Step 5 --- leaf direction constraints (\file{:204--213}).} Each
external leaf is classified by the edges touching it: an only-outgoing leaf can
carry only a response field (\code{'resp'}); only-incoming, only a physical field
(\code{'phys'}); a leaf with both is \code{'both'}.

\paragraph{Step 6 --- external-field permutations (\file{:223--242}).} This is
the step a reader is most likely to find surprising, so it gets its own
subsection below. The enumerator loops over \emph{every distinct ordering} of
\code{external\_fields} across the leaves, keeping only orderings in which each
leaf's field is compatible with its direction.

\paragraph{Step 7 --- dispatch (\file{:244--266}).} If there are no internal
vertices, the edges merely connect leaves, and the work is delegated to
\code{\_try\_build\_diagram\_no\_internal}. Otherwise the enumerator takes the
Cartesian product \code{product(*candidate\_lists)} over the internal vertices'
candidate type lists, and for each combination builds a \code{vert\_assignment}
and delegates to \code{\_try\_build\_diagram}.

\subsection{Why enumerate external-field permutations at all?}

Step 6 deserves its own explanation because it looks like double work. The
relevant code is:

\begin{lstlisting}
for ext_perm in _distinct_permutations(tuple(external_fields)):  # type_assignment.py:223
    ext_assignment = {}
    valid_ext = True
    for leaf_idx in range(len(ext_perm)):
        lf = leaves[leaf_idx]
        field = ext_perm[leaf_idx]
        direction = leaf_directions[lf]
        if direction == 'resp' and field not in resp_index:
            valid_ext = False; break
        if direction == 'phys' and field not in phys_index:
            valid_ext = False; break
        if direction == 'both':
            if field not in resp_index and field not in phys_index:
                valid_ext = False; break
        ext_assignment[lf] = field
    if not valid_ext:
        continue
    ...
\end{lstlisting}

The upstream prediagram dedup (Phase~C) treats all leaves as interchangeable
stubs, so a \emph{single} prediagram stands in for every way of distributing the
external fields across its leaves. To recover the diagrams that differ in
\emph{which leaf carries which field} --- for example \code{dn2} sitting on a
source-leaf versus on an interaction-leaf --- we must materialise those
distributions here. The helper \code{\_distinct\_permutations}
(\file{type\_assignment.py:22--24}) is a one-liner that returns
\code{set(permutations(seq))}, so two identical external fields do not produce two
``different'' orderings. Any orderings that nevertheless turn out to be truly
identical diagrams are merged later by Phase~G's deduplication.

\subsection{The degenerate case: a bare propagator}

When a prediagram has \emph{no} internal vertices --- e.g. a single edge directly
joining two external leaves, which is just a bare propagator ---
\code{\_try\_build\_diagram\_no\_internal} (\file{type\_assignment.py:269--305})
handles it directly. For each edge $u\to v$: the tail $u$ contributes the
response leg (its external field), the head $v$ the physical leg; both must be
direction-compatible; we look up \code{ri = resp\_index[resp\_field]} and
\code{pi = phys\_index[phys\_field]}; and we reject the edge if
\code{g\_zero\_mask[(pi, ri)]} is true --- note again the \textbf{(pi, ri)} order.
On success it records the edge's types and propagator indices and yields a
\code{TypedDiagram} with empty \code{vertex\_assignments}.

\begin{example}[A single edge, end to end]
The micro-example from \file{tests/test\_type\_assignment.py} makes the data
shapes concrete. With a prediagram that is one directed edge $1\to2$ between two
leaves, and external fields \code{[('nt', 1), ('dn', 1)]}:
\begin{lstlisting}
pd = (DiGraph([(1, 2)]), None, [1, 2], [])
external_fields = [('nt', 1), ('dn', 1)]
results = list(enumerate_typed_diagrams(pd, external_fields, [], [],
                                        G_ft, resp_idx, phys_idx))
# -> one TypedDiagram with
#    edge_types         = {(1, 2, None): (('nt',1), ('dn',1))}
#    external_legs      = {1: ('nt',1), 2: ('dn',1)}
#    propagator_indices = {(1, 2, None): (ri, pi)}
#    vertex_assignments = {}
\end{lstlisting}
The response field \code{nt} sits at the tail, the physical field \code{dn} at
the head, and the edge carries $G_{\mathrm{ft}}[pi, ri]$ --- one bare propagator,
exactly as Fig.~\ref{fig:ta-edge} drew it.
\end{example}

\subsection{Leg matchings and the backtracker}

For a prediagram \emph{with} internal vertices, \code{\_try\_build\_diagram}
(\file{type\_assignment.py:308--337}) sets up the search. For each internal
vertex it computes the list of ways to map that vertex type's legs onto its
incident edges, via \code{\_leg\_matchings}; if any vertex admits zero such
matchings the type combination is dead and it \code{return}s. Otherwise it
launches the recursive backtracker.

\code{\_leg\_matchings(vertex\_type, out\_edges, in\_edges)}
(\file{type\_assignment.py:340--416}) enumerates the distinct bijections from a
vertex's legs to its edges: each outgoing edge must receive a \emph{response}
leg, each incoming edge a \emph{physical} leg. Its core is short:

\begin{lstlisting}
resp_legs = vertex_type.response_legs
has_phys = hasattr(vertex_type, 'physical_legs')
phys_legs = vertex_type.physical_legs if has_phys else []     # sources have none
if len(resp_legs) != len(out_edges) or len(phys_legs) != len(in_edges):
    return
from sympy.utilities.iterables import multiset_permutations    # type_assignment.py:406
resp_perms = (list(multiset_permutations(list(resp_legs))) if resp_legs else [[]])
phys_perms = (list(multiset_permutations(list(phys_legs))) if phys_legs else [[]])
for rp in resp_perms:
    resp_map = dict(zip(out_edges, rp))
    for pp in phys_perms:
        phys_map = dict(zip(in_edges, pp))
        yield resp_map, phys_map
\end{lstlisting}

The \code{has\_phys} check is how a source is told apart from an interaction
vertex: a \code{SourceType} has no \code{physical\_legs} attribute at all, so
\code{phys\_legs} defaults to the empty list and the source contributes only
response legs. The choice of \code{multiset\_permutations} over a full
\code{N!} enumeration is the single most important optimisation in the file, and
\S\ref{sec:ta-canonical} is devoted to proving it is correct.

The recursive core \code{\_backtrack} (\file{type\_assignment.py:419--538})
carries three pieces of state: \code{vertex\_idx} (which internal vertex we are
assigning), and two partial maps \code{assigned\_resp} and \code{assigned\_phys}
(\code{\{edge: leg\}} so far). It has two parts.

\paragraph{The base case (\file{:431--472}).} All internal vertices are assigned.
For each edge, the code resolves its response leg --- from \code{assigned\_resp}
if an internal tail set it, else from the external field if the tail is a leaf
--- and, symmetrically, its physical leg. It looks up \code{ri}/\code{pi}; if
either is missing or the propagator is zero (\code{g\_zero\_mask[(pi, ri)]}), it
\code{return}s. On full success it yields the finished \code{TypedDiagram}.

\paragraph{The recursive step (\file{:474--538}).} For each leg-matching option at
the current vertex, it merges that option into fresh \code{new\_resp}/
\code{new\_phys} copies and runs an \emph{early propagator-consistency check} on
every edge that is now fully determined --- both ends known. It checks both
internal-internal edges and edges where one end is a leaf, and prunes immediately
if any propagator is zero. Only if the partial assignment is still
\code{consistent} does it recurse to \code{vertex\_idx + 1}. This early pruning
is what keeps the search tractable on large diagrams: a dead branch is abandoned
the moment a single edge is forced onto a zero propagator, rather than after a
full assignment.

\begin{gotcha}
The convention note at \file{type\_assignment.py:490--502} is worth reading in
the source. It states explicitly that the correct lookup order is
\code{g\_zero\_mask[(pi, ri)]} (physical-row, response-col) and records a real
historical bug: a previous version used \code{[ri, pi]}, which happened to read
the right entry for \emph{diagonal} couplings (where the response and physical
indices coincide) but silently filtered out valid diagrams for \emph{off-diagonal}
ones. The cited example is a GTaS coupling $\tilde m \to \delta n$ with physical
index $0$ and response index $4$: there \code{[ri, pi]} and \code{[pi, ri]} point
at different cells, so the wrong order rejected diagrams whose propagator was in
fact nonzero. Any new propagator-touching code must index $(pi, ri)$.
\end{gotcha}

% ---------------------------------------------------------------------
\subsection{The canonical-leg-matching optimisation}\label{sec:ta-canonical}

Why does \code{\_leg\_matchings} use SymPy's \code{multiset\_permutations}
instead of \code{itertools.permutations}? The difference is a factor that can be
$15$--$20\times$ on the size of the intermediate diagram set, and getting it
right without changing the final number is a small theorem the code's docstring
(\file{type\_assignment.py:356--387}) carefully proves.

\begin{defn}
\textbf{Multiset permutations.}
For a multiset of length $N$ with multiplicities $n_r$ (i.e. $n_r$ copies of the
$r$-th distinct element), the number of \emph{distinct} orderings is the
multinomial coefficient
\[
  \frac{N!}{\prod_r n_r!},
\]
not the full $N!$. SymPy's \code{multiset\_permutations(seq)} yields each distinct
ordering exactly once. Because the leg values are hashable \code{(field\_base,
pop\_idx)} tuples, the multiset comparison just works.
\end{defn}

The argument for correctness is this. Earlier code enumerated all $N!$ index
permutations and \emph{relied on} the downstream deduplication to discard the
overcount. That overcount factor is \emph{exactly} the symmetry factor
$\Sym(\Gamma)$. So the old pipeline generated $\Sym(\Gamma)$ identical copies of
each diagram, each copy carrying its full weight; dedup kept one; and the
integrator then multiplied by $\Sym(\Gamma)$. The new code skips the inflation
entirely: it generates the single \emph{canonical} ordering directly, and the
integrator multiplies by the \emph{same} $\Sym(\Gamma)$. The output is
bit-identical; only the intermediate count (and the cold-start enumeration wall
time) shrinks. This equivalence is pinned by the regression tests
\code{test\_leg\_matchings\_canonical} and
\code{test\_enumerate\_typed\_signatures\_match\_pre\_change} in
\file{tests/test\_type\_assignment.py}.

\begin{gotcha}
This optimisation is correct \emph{only because} the integrator multiplies by
$\Sym(\Gamma)$ afterward. The two are a matched pair. If you ever change how the
symmetry factor is applied downstream, you must revisit \code{\_leg\_matchings} ---
the regression tests above exist precisely to catch a drift between the two.
\end{gotcha}

% ---------------------------------------------------------------------
\subsection{The propagator zero mask, and why a string test}\label{sec:ta-zeromask}

We deferred one design choice from Step 1: why is the zero test
\code{str(G\_ft[i,j]) == '0'} rather than the obviously-correct
\code{G\_ft[i,j].is\_zero()}? The answer is entirely about cost, and the code's
comment (\file{type\_assignment.py:144--157}) spells it out.

$G_{\mathrm{ft}}$ entries can be long rational functions of the symbolic
parameters. Calling \code{is\_zero()} (or, equivalently, \code{e == SR(0)}) on
such an entry triggers Sage's full symbolic canonicalisation engine, which can
cost \emph{several seconds per entry} on a rich theory --- and the backtracker
would otherwise hit that cost on every leg assignment. The string test sidesteps
canonicalisation completely: it returns \code{True} only for entries that are
\emph{literally} the symbol \code{SR(0)}, which is what the propagator builder
writes into empty cells of $K^{-1}$.

The trade-off is a deliberate, harmless over-generation. An entry that is zero
only \emph{after cancellation} --- a nontrivial expression that happens to
simplify to $0$ --- is \emph{not} caught by the string test, so it slips through
and a diagram or two carrying it gets generated. But those diagrams simply
integrate to $0$ downstream. The string test can therefore over-generate, never
under-generate: it never rejects a valid diagram, so it can never produce a wrong
number, only a few extra ``vanishing'' diagrams.

\begin{note}
This is a recurring pattern in the codebase: when an exact symbolic decision is
expensive and a cheap structural proxy errs only on the safe side (here,
over-generation that integrates to zero), the cheap proxy wins. Be aware, though,
that the typed-diagram set you inspect may contain a handful of diagrams that
contribute nothing --- this is by design, not a bug.
\end{note}

% =====================================================================
\section{Phase F --- causality}\label{sec:ta-causality}

Phase F (\file{msrjd/diagrams/causality.py}) prunes typed diagrams that violate
the retarded boundary condition. It is short, because in the current pipeline
only the structural DAG check is live.

\subsection{The per-diagram check}

\begin{lstlisting}
def check_causality(typed_diagram, pole_vals=None):  # causality.py:97-123
    D = typed_diagram.prediagram[0]
    # Structural check: the prediagram must be a DAG
    if not D.is_directed_acyclic():
        return False, 'Prediagram contains a directed cycle - acausal.', []
    # Pole check (if available)
    if pole_vals is not None and len(pole_vals) > 0:
        return check_pole_structure(pole_vals)
    return True, 'Structural DAG check passed. No poles to verify.', []
\end{lstlisting}

It pulls the directed graph \code{D = typed\_diagram.prediagram[0]} and asks
\code{D.is\_directed\_acyclic()}. A directed cycle means you can return to a
vertex by following retarded arrows --- a closed loop of strictly-retarded
propagators --- which is acausal, so the diagram is rejected. Only if
\code{pole\_vals} are supplied does it fall through to the analytic pole check;
in the live pipeline none are passed, so \textbf{the DAG check is the operative
gate}.

\code{filter\_causal(typed\_diagrams)} (\file{causality.py:126--152}) is the
list-level wrapper the orchestrator calls (\file{pipeline/\_diagrams.py:185}): it
runs \code{check\_causality} on each diagram and partitions them into a
\code{kept} list and a list of one human-readable rejection reason per discarded
diagram, returning \code{(kept, n\_discarded, discarded\_details)}.

\subsection{The latent pole machinery}

\code{check\_pole\_structure(pole\_vals, omega=None)}
(\file{causality.py:19--94}) is the analytic counterpart, present but dormant.
Given a list of pole locations it coerces each with \code{SR(pole)}, takes its
imaginary part with Sage's \code{imag\_part}, attempts \code{simplify\_full()},
and then tries to decide the sign:

\begin{itemize}
  \item \code{bool(im\_simplified > 0)} --- the pole passes (upper half-plane);
  \item \code{bool(im\_simplified <= 0)} --- the pole fails;
  \item \code{bool(im\_simplified == 0)} --- a pole exactly on the real axis is
        treated as a \emph{failure} (it is not retarded).
\end{itemize}

Each \code{bool(...)} is wrapped in \code{try/except (TypeError, ValueError)}
because Sage raises when it cannot decide a symbolic inequality; an undecidable
pole is recorded as \emph{conditional}. The final verdict: any failure
$\Rightarrow$ \code{False}; else any conditional $\Rightarrow$ \code{True}
\emph{with} the symbolic $\operatorname{Im}(\text{pole}) > 0$ constraints returned
as \code{conditions}; else everything passed. So an inconclusive symbolic pole
returns \code{passed=True} (carrying its constraint), while a real-axis pole is a
hard fail.

% =====================================================================
\section{Phase G --- the symmetry factor and deduplication}\label{sec:ta-symmetry}

Phase G (\file{msrjd/diagrams/symmetry.py}) does the two jobs whose
correctness matters most, because a mistake here produces a numerically
\emph{plausible} but \emph{wrong} answer rather than an obvious crash. First, it
deduplicates the typed diagrams. Second, it computes the symmetry factor
$\Sym(\Gamma)$. Both rest on a single construction --- a coloured incidence
digraph fed to \code{nauty} --- so we build that first.

\subsection{The canonical Feynman rule}

The code implements what its comments call ``Path A'': the canonical Feynman
rule (\file{symmetry.py:403--446}),
\begin{equation}
  \Sym(\Gamma) \;=\;
  \frac{\displaystyle\prod_v \prod_\ell n_{v,\ell}!}
       {\bigl|\Aut_{\text{fixed ext}}(\Gamma)\bigr|}.
  \label{eq:ta-sym}
\end{equation}
There are two pieces, a numerator and a denominator, and the per-diagram weight
the integrator ultimately builds (\file{symmetry.py:421--423}) is
\begin{equation}
  \mathrm{weight}(\Gamma) \;=\;
  \Sym(\Gamma)\;\times\;\prod_v c_{\text{user},v}\;\times\;\prod_e P_e
  \;\times\;\int \dd t_v,
  \label{eq:ta-weight}
\end{equation}
where $c_{\text{user},v}$ is the \emph{literal} coefficient written in the action
(no implicit $1/n!$), $P_e$ is the propagator on edge $e$, and the integral runs
over the internal vertex times.

\paragraph{The numerator --- the Wick leg factor.} For each vertex $v$ and each
leg-type $\ell$ (a leg-type lumps \emph{direction} --- response vs physical ---
with the \emph{field}), let $n_{v,\ell}$ be the number of legs of type $\ell$ at
$v$. The numerator is the product of $n_{v,\ell}!$ over all of them --- ``the Wick
combinatorial you would get if every leg at every vertex were distinguishable.''
A degree-3 vertex $\varphi^2\tilde\varphi$ contributes $2!$ from its two
identical $\varphi$ legs. This is \code{\_wick\_leg\_factor}
(\file{symmetry.py:115--146}):

\begin{lstlisting}
def _wick_leg_factor(typed_diagram):  # symmetry.py:115-146
    n = 1
    for v, vtype in typed_diagram.vertex_assignments.items():
        resp_counts = Counter(vtype.response_legs)      # outgoing/response legs
        for c in resp_counts.values():
            n *= factorial(c)
        if hasattr(vtype, 'physical_legs'):             # incoming/physical legs
            phys_counts = Counter(vtype.physical_legs)
            for c in phys_counts.values():
                n *= factorial(c)
    return n
\end{lstlisting}

\code{Counter} (from the standard library's \code{collections}) tallies the
multiplicity of each distinct leg, and \code{factorial} (from \code{math}) is
$n!$.

\paragraph{The denominator --- the automorphism order.}
$|\Aut_{\text{fixed ext}}(\Gamma)|$ is the order of the automorphism group of the
\emph{fully coloured} diagram, with the external legs \emph{pinned} at their
canonical positions. An \term{automorphism} is a relabelling of internal vertices
and edges that maps the diagram to itself while preserving every label physics
cares about: vertex type, vertex coefficient, the propagator on each edge, and
the direction of each edge. It captures three families of redundancy ---
same-type internal vertex swaps, parallel-edge swaps between the same vertex
pair, and self-loop / tadpole leg-pair swaps --- and dividing the Wick numerator
by it removes exactly the over-counting those redundancies introduce.

\subsection{The coloured incidence digraph}

The obstacle is that \code{nauty} (via Sage) colours \emph{vertices}, not edges
--- yet an automorphism must respect the colour (the propagator type) of each
\emph{edge}. The resolution, in \code{\_colored\_incidence\_digraph}
(\file{symmetry.py:149--244}), is to build a fresh bipartite ``vertex-nodes plus
edge-nodes'' graph:

\begin{itemize}
  \item each Feynman vertex becomes a node \code{('V', v)};
  \item \emph{each propagator edge becomes its own node} \code{('E', ek)};
  \item an original edge $u\to v$ is encoded as \emph{two} bipartite edges,
        \code{('V', u) -> ('E', ek)} and \code{('E', ek) -> ('V', v)}.
\end{itemize}

Now an edge's colour lives on its edge-node, where \code{nauty} can see it. The
node colours (the \emph{partition}, in nauty's language) are:

\begin{description}
  \item[Internal vertex node] colour
        \code{('vertex', type(vt).\_\_name\_\_, str(vt.coefficient), vt.bigrade)}
        --- so only same-type, same-coefficient, same-bigrade vertices may be
        swapped.
  \item[External leaf node] if \code{fix\_external=True}, colour
        \code{('leaf', v, field)} --- each leaf gets its own colour, so leaves can
        never move; this gives $\Aut_{\text{fixed ext}}$. If
        \code{fix\_external=False}, colour \code{('leaf', field)} --- same-field
        leaves share a colour and so \emph{may} permute; this gives
        $\Aut_{\text{free ext}}$.
  \item[Edge node] colour \code{('edge', resp\_leg, phys\_leg, prop)}, where
        \code{prop} is the $(ri, pi)$ propagator-index tuple --- so an
        automorphism must preserve the propagator type \emph{and} direction on
        every edge.
\end{description}

Parallel edges between the same vertex pair with identical colour become
genuinely swappable edge-nodes, which is exactly the symmetry the Feynman rules
want counted. The function returns the digraph \code{D}, the \code{partition}
(the list of colour classes, ready to hand to nauty), and bookkeeping maps.

\subsection{Counting automorphisms with nauty}

With the coloured digraph in hand, \code{\_automorphism\_order}
(\file{symmetry.py:328--340}) is the single place the automorphism count is
invoked:

\begin{lstlisting}
def _automorphism_order(typed_diagram, fix_external=True):  # symmetry.py:328-340
    D, partition, _, _ = _colored_incidence_digraph(
        typed_diagram, fix_external=fix_external)
    grp = D.automorphism_group(partition=partition)
    return int(grp.order())
\end{lstlisting}

\code{D.automorphism\_group(partition=partition)} asks Sage --- and underneath,
\code{nauty} --- for the group of colour-preserving relabellings; \code{grp} is a
Sage permutation group, and we keep only its \code{order()} (an integer),
discarding the group itself. Passing \code{fix\_external=True} yields
$|\Aut_{\text{fixed ext}}|$; \code{False} yields $|\Aut_{\text{free ext}}|$.

The production symmetry factor \code{combinatorial\_factor}
(\file{symmetry.py:403--446}) then divides numerator by denominator:

\begin{lstlisting}
def combinatorial_factor(typed_diagram):  # symmetry.py:403-446
    numer = _wick_leg_factor(typed_diagram)
    aut = _automorphism_order(typed_diagram, fix_external=True)
    if aut <= 0 or numer % aut != 0:
        return numer // max(aut, 1)     # defensive floor-division fallback
    return numer // aut
\end{lstlisting}

The integer division should always be exact --- a finite group acting freely on a
finite set --- and the floor-division branch is purely defensive against a Sage
edge case. \code{compute\_all\_combinatorial\_factors}
(\file{symmetry.py:449--462}) is the obvious list comprehension over this.

\begin{note}
\file{symmetry.py} also carries an older, vertex-local way to count Wick
pairings: \code{\_vertex\_combinatorial\_factor} (\file{symmetry.py:51--112}),
which computes a per-vertex factor $M_v$ and would give
$\Sym(\Gamma) = \prod_v M_v$ in the ``Path B'' formulation. The production
\code{combinatorial\_factor} does \emph{not} use it --- it uses the Aut-based
Eq.~\eqref{eq:ta-sym}. The $M_v$ function is retained because it is the
per-vertex semantic that \S\ref{sec:ta-canonical}'s leg-matching optimisation
appeals to (the overcount factor it avoids generating is exactly this $M_v$
product), and as a documentation cross-check.
\end{note}

% ---------------------------------------------------------------------
\subsection{Deduplication and the complete invariant}\label{sec:ta-dedup}

Phase E over-generates: it emits the same diagram many times, relabelled and
leaf-permuted. Deduplication collapses those copies, and the key it uses must be
a \emph{complete} graph-isomorphism invariant --- two typed diagrams must hash
equal \emph{if and only if} they are isomorphic coloured graphs. A weaker
invariant that occasionally collides non-isomorphic diagrams would silently
\emph{drop} one class's integral, which is the worst kind of bug: a wrong number
with no error.

\code{diagram\_signature(td)} (\file{symmetry.py:467--524}) provides exactly such
a complete invariant, built from \code{nauty}'s canonical labelling:

\begin{lstlisting}
def diagram_signature(td):  # symmetry.py:467-524
    D, _, _, color_groups = _colored_incidence_digraph(td, fix_external=False)
    keys = sorted(color_groups.keys(), key=str)         # label-independent order
    partition = [color_groups[k] for k in keys]
    C, cert = D.canonical_label(partition=partition, certificate=True)
    cells = tuple(tuple(sorted(cert[v] for v in color_groups[k])) for k in keys)
    edges = tuple(sorted(C.edges(labels=False)))
    return (tuple(str(k) for k in keys), cells, edges)
\end{lstlisting}

\begin{defn}
\textbf{Canonical label and certificate.}
\code{nauty} computes a \term{canonical form}: a relabelling of a graph's
vertices into a normal order such that two graphs are isomorphic \emph{iff} their
canonical forms are identical. Sage exposes it as
\code{D.canonical\_label(partition=..., certificate=True)}, which returns a pair
\code{(C, cert)}: \code{C} is the canonically relabelled copy of \code{D}, and the
\term{certificate} \code{cert} is the dict mapping each original vertex id to its
canonical-label id.
\end{defn}

The signature is built from three pieces: the colour keys (as strings, for
hashability), \emph{which canonical ids each colour class occupies}
(\code{cells}), and the canonical edge list (\code{edges}). The \code{cells} piece
is essential --- the canonical edge list \emph{alone} would confuse two graphs of
the same shape but different colourings, so we must also record where each colour
landed. It is built with \code{fix\_external=False} so that same-field external
leaves are allowed to permute; leaf-permuted variants are thereby merged here and
re-expanded downstream by the integrator's mapping sum (\S\ref{sec:ta-comp}).

\code{deduplicate\_with\_multiplicities(typed\_diagrams)}
(\file{symmetry.py:553--593}) is the function the orchestrator actually calls
(\file{pipeline/\_diagrams.py:186}). It keeps one representative per signature in
first-seen order, alongside each equivalence class's size:

\begin{lstlisting}
sig_to_idx = {}
unique = []
multiplicities = []
for td in typed_diagrams:
    sig = diagram_signature(td)
    idx = sig_to_idx.get(sig)
    if idx is None:
        sig_to_idx[sig] = len(unique)
        unique.append(td)
        multiplicities.append(1)
    else:
        multiplicities[idx] += 1
return unique, multiplicities
\end{lstlisting}

(\code{deduplicate\_typed\_diagrams} (\file{symmetry.py:527--550}) is a thin
wrapper returning just the \code{unique} list.)

\begin{gotcha}
The \code{multiplicities} are \textbf{diagnostic only} --- never multiply them
back into the weight. Under Path A the entire weight is carried by
\code{combinatorial\_factor}; multiplying by the class size would double-count,
because the merged entries are isomorphic copies whose pairings $\Sym(\Gamma)$
already includes (\file{symmetry.py:556--568} is explicit). This subtlety has
teeth: the on-disk cache stage tag was bumped (to \code{unique\_typed\_mult\_v3})
specifically to invalidate caches predating the complete-invariant signature,
since loading them would resurrect the old class-collision bug.
\end{gotcha}

\begin{example}[The collision bug the complete invariant closed]
The history note at \file{symmetry.py:484--498} records the concrete failure. At
$k=3$, one loop, the OU$+ax^2+bx^3$ theory's $a^3$ sector has $11$ distinct
diagram classes. A previous, hand-rolled signature --- a depth-1 isomorphism
\emph{invariant} rather than a complete one --- collided $4$ of those $11$ into
$2$ representatives. Because the dedup multiplicity is correctly \emph{not}
multiplied back, the collided classes' integrals were silently dropped, and the
$\kappa_3$ one-loop coefficient came out $-\tfrac{68}{3}a^3$ instead of the
correct $-32\,a^3$. A canonical-form invariant \emph{cannot} collide
non-isomorphic graphs, so this failure mode is closed for every $(k, \ell,
\text{theory})$. The fix was validated against the exact Boltzmann series and a
brute-force labelled Wick enumeration per class.
\end{example}

% ---------------------------------------------------------------------
\subsection{External Wick compensation and the integrator hand-off}\label{sec:ta-comp}

There are \emph{two} automorphism orders, differing only in whether external
leaves may permute, and their ratio is the bridge to the integrator. With leaves
pinned we get $|\Aut_{\text{fixed ext}}|$ (the denominator of $\Sym(\Gamma)$);
with same-field leaves free we get $|\Aut_{\text{free ext}}|$. Since
$\Aut_{\text{fixed ext}}$ is a subgroup of $\Aut_{\text{free ext}}$, the latter
order is always an integer multiple of the former, and that integer is the
\term{external Wick compensation}:
\begin{equation}
  \mathrm{comp} \;=\;
  \frac{|\Aut_{\text{free ext}}(\Gamma)|}{|\Aut_{\text{fixed ext}}(\Gamma)|}.
  \label{eq:ta-comp}
\end{equation}
This is \code{external\_wick\_compensation} (\file{symmetry.py:343--400}):

\begin{lstlisting}
def external_wick_compensation(typed_diagram):  # symmetry.py:343-400
    aut_free = _automorphism_order(typed_diagram, fix_external=False)
    aut_fixed = _automorphism_order(typed_diagram, fix_external=True)
    if aut_fixed <= 0 or aut_free % aut_fixed != 0:
        raise RuntimeError(
            f'external_wick_compensation: |Aut_free|={aut_free} not '
            f'divisible by |Aut_fixed|={aut_fixed}')
    return aut_free // aut_fixed
\end{lstlisting}

Why it exists: the integrators enumerate \emph{every} field-respecting assignment
of canonical external positions to a diagram's leaves and sum the integrand over
all of them. By the \term{orbit--stabilizer theorem} ($|\text{orbit}| =
|\text{group}| / |\text{stabilizer}|$), two assignments give the same
pinned-external diagram iff they differ by a leaf-permuting automorphism, and the
stabilizer of any assignment is exactly $\Aut_{\text{fixed ext}}$. So each
distinct pinned diagram is counted $|\text{free}|/|\text{fixed}|$ times by the
mapping sum; dividing by \code{comp} makes the sum equal the sum over distinct
pinned diagrams --- which is what the Feynman rules require, since each pinned
diagram already carries its full weight through $\Sym(\Gamma)$, whose denominator
is the \emph{same} $\Aut_{\text{fixed ext}}$.

\begin{gotcha}
\code{external\_wick\_compensation} fails \emph{loud}: a non-divisor raises
\code{RuntimeError} rather than guessing (\file{symmetry.py:391--399}). This is a
deliberate asymmetry with \code{combinatorial\_factor}, which falls back to
floor-division (\file{symmetry.py:440--445}). The reasoning: a wrong compensation
silently \emph{mis-weights} diagrams --- the exact bug class this function exists
to fix --- whereas the Wick factor's divisibility is a theorem that should always
hold. So the compensation refuses to guess, while the Wick factor is allowed a
defensive fallback.
\end{gotcha}

\begin{gotcha}
A nearby function, \code{vertex\_role\_signature} (\file{symmetry.py:247--325}),
is a hashable per-vertex signature invariant under relabelling. It once served as
a Phase~J external divisor (a $\prod N!$ over leaves grouped by this signature),
but that depth-1 heuristic \emph{over-divided} when two structurally distinct
vertices happened to share a signature without any automorphism relating them ---
breaking every $k\ge3$ one-loop cumulant. It has been \emph{superseded} by the
exact \code{external\_wick\_compensation} of Eq.~\eqref{eq:ta-comp}. The function
survives in the file but the production weight path no longer uses it; do not
reintroduce the heuristic divisor.
\end{gotcha}

\subsection{Coefficient classification}\label{sec:ta-coeff}

One more Phase G function is called directly by the integrators:
\code{classify\_coefficient\_factors} (\file{symmetry.py:618--698}). Its job is to
split each vertex coefficient into a \emph{scalar prefactor} that can be pulled
outside the integral versus \emph{factors that must stay inside} (time-dependent
couplings, non-white noise amplitudes). It returns a dict with keys \code{'Scal'}
(the symmetry factor $\Sym(\Gamma)$ itself), \code{'scalar\_prefactor'} (the
product of all pull-out-able pieces, including $\Sym(\Gamma)$),
\code{'vertex\_time\_factors'}, \code{'source\_time\_info'}, and
\code{'is\_stationary'}.

The mechanics, at a high level: it seeds \code{scalar\_parts} with \code{SR(Scal)}
and, for each vertex, takes \code{coeff = -SR(vtype.coefficient)} --- the minus
sign coming from $Z = \int \ee^{-S}$ in Eq.~\eqref{eq:ta-Z}, so each interaction
factor from $S$ acquires a sign. For a \emph{source}, white noise with a
time-independent amplitude is pulled out; coloured or time-dependent amplitudes
stay inside and are recorded in \code{source\_time\_info}. For an
\emph{interaction vertex}, the constant part (obtained by substituting every
time-dependent symbol to $1$) joins \code{scalar\_parts}, and the residual
\code{td\_part = coeff / const\_part} (tidied with \code{simplify\_rational})
goes into \code{vertex\_time\_factors}. Finally
\code{scalar\_prefactor = reduce(mul, scalar\_parts, SR(1))} multiplies the
pull-out pieces into one product, and \code{is\_stationary} is \code{True} when
there are no vertex time factors and no time-dependent or non-standard noise
amplitudes.

% =====================================================================
\section{The orchestrator: wiring the four phases}\label{sec:ta-orchestrator}

All four phases are wired together --- with disk caching --- by
\code{enumerate\_unique\_diagrams} in \file{pipeline/\_diagrams.py}. The heart of
it, after a cache miss, is four lines
(\file{pipeline/\_diagrams.py:176--186}):

\begin{lstlisting}
_, _, prediagrams, _ = enumerate_prediagrams_all(k=k, ell=ell, verbose=False)
typed = enumerate_all_typed(prediagrams, external_fields, vtypes, stypes,
                            G_ft=G_ft, resp_index=resp_idx, phys_index=phys_idx,
                            parallel=parallel, n_workers=n_workers)
causal, n_disc, _ = filter_causal(typed)
unique, multiplicities = deduplicate_with_multiplicities(causal)
\end{lstlisting}

Read top to bottom this is the whole chapter in miniature: Phase~C produces the
prediagrams; \code{enumerate\_all\_typed} (the cross-prediagram driver of Phase~E,
which we meet by its full name \code{enumerate\_all} in \S\ref{sec:ta-fork})
labels them; \code{filter\_causal} (Phase~F) keeps the DAGs; and
\code{deduplicate\_with\_multiplicities} (Phase~G) collapses isomorphic copies.
The flow of intermediate objects is:

\begin{center}
\begin{tikzpicture}[
    >={Stealth[length=2.4mm]},
    box/.style={draw,rounded corners,fill=codebg,align=center,
                font=\footnotesize,inner sep=4pt,minimum height=9mm},
    lbl/.style={font=\scriptsize\itshape,midway,above},
  ]
  \node[box] (pc)  at (0,0)    {prediagrams\\\scriptsize $(D,G,\text{leaves},\text{internal})$};
  \node[box] (pe)  at (3.9,0)  {typed\\\scriptsize list[TypedDiagram]};
  \node[box] (pf)  at (7.5,0)  {causal\\\scriptsize list[TypedDiagram]};
  \node[box] (pg)  at (11.1,0) {unique\\\scriptsize $+$ multiplicities};
  \draw[->] (pc) -- (pe) node[lbl] {E};
  \draw[->] (pe) -- (pf) node[lbl] {F};
  \draw[->] (pf) -- (pg) node[lbl] {G};
\end{tikzpicture}
\end{center}

Note that Phase~D does not appear: as \S\ref{sec:ta-filter} explained, the loop
enumerator already emits degree-valid prediagrams, so the explicit filter is
skipped on the live path.

Downstream, the integrators (Phase~J) take each unique \code{TypedDiagram} and,
per diagram: read \code{vertex\_assignments} for coefficients (split via
\code{classify\_coefficient\_factors}), read \code{edge\_types} and
\code{propagator\_indices} to place propagators $G_{\mathrm{ft}}[pi, ri]$ on
edges, read \code{external\_legs} to line the leaves up with
\code{external\_fields}, multiply by $\Sym(\Gamma) =
\code{combinatorial\_factor(td)}$, sum over external-leaf mappings while dividing
by \code{external\_wick\_compensation(td)}, and finally perform the time /
frequency / momentum integral.

% =====================================================================
\section{Parallelism and fork-safety}\label{sec:ta-fork}

The cross-prediagram driver of Phase~E, \code{enumerate\_all}
(\file{type\_assignment.py:582--684}), is the only function here with a
performance story worth telling, and it comes with a genuine landmine. Labelling
one prediagram is independent of labelling any other --- the work is
\emph{embarrassingly parallel} across prediagrams --- so \code{enumerate\_all}
offers an optional fork-based multiprocessing path. The serial and parallel paths
produce \textbf{bit-identical} output, in the same order, pinned by
\code{test\_enumerate\_all\_parallel\_matches\_serial}.

\subsection{Why fork, and why it is lethal in a notebook}

The path uses POSIX \code{fork} specifically because Sage \code{DiGraph} objects,
\code{VertexType}/\code{SourceType} instances, and the propagator matrix do not
all pickle cleanly through stdlib \code{pickle} --- so the ``spawn'' start method,
which passes inputs to workers by pickling, would fail. With \code{fork}, the
workers simply \emph{inherit} the parent's memory and already have those objects;
only the \emph{return} value (a list of \code{TypedDiagram}, which does pickle via
its \code{\_\_getstate\_\_}) crosses back. The inputs are stashed in a module
global \code{\_ENUM\_WORKER\_STATE} \emph{before} the fork so the children inherit
them rather than receiving them by pickling.

The catastrophe is platform-specific. On \textbf{macOS}, Apple's runtime forbids
using many Objective-C / Cocoa APIs after a \code{fork()} that is not followed by
an \code{exec()}. Once Cocoa, BLAS, or matplotlib has initialised in the parent
--- which \emph{always} happens inside a Jupyter kernel --- forking and then
touching those libraries in the child can \textbf{hard-crash the kernel and the
entire OS}. (This crashed the project's dev machine more than once.) The
``fix'' people reach for, setting the environment variable
\code{OBJC\_DISABLE\_INITIALIZE\_FORK\_SAFETY=YES}, makes it \emph{worse}: it
converts a controlled abort into a hard crash.

\subsection{The guard}

The shared guard lives in \file{msrjd/fork\_safety.py} and is consulted by every
fork-based path in the repository (Phase~E here, and the temporal Phase~J
batch driver). Its core predicate returns \code{True} only in the precise lethal
situation:

\begin{lstlisting}
def fork_unsafe_in_notebook(start_method='fork'):  # fork_safety.py:27-44
    if start_method != 'fork':
        return False
    if sys.platform != 'darwin':
        return False
    try:
        from IPython import get_ipython
        ip = get_ipython()
    except Exception:
        return False
    return ip is not None and ip.__class__.__name__ == 'ZMQInteractiveShell'
\end{lstlisting}

The Jupyter fingerprint is the class-name check \code{ZMQInteractiveShell} --- the
IPython shell class that \emph{only} a ZMQ/Jupyter kernel uses. Terminal IPython
uses \code{TerminalInteractiveShell}; plain scripts and pytest return \code{None}
from \code{get\_ipython()}. So fork stays available --- and fast --- everywhere it
is safe (scripts run as \code{sage -python run.py}, pytest, Linux, terminal
IPython), and is suppressed only inside a macOS Jupyter kernel.

When the guard fires, \code{enumerate\_all} warns once and degrades to serial
(\file{type\_assignment.py:630--634}):

\begin{lstlisting}
from msrjd.fork_safety import fork_unsafe_in_notebook, warn_fork_fallback_once
if (parallel and len(prediagrams) > 1
        and fork_unsafe_in_notebook(start_method)):
    warn_fork_fallback_once('diagram type-assignment enumeration')
    parallel = False
\end{lstlisting}

\code{warn\_fork\_fallback\_once(where)} (\file{fork\_safety.py:47--61}) emits a
single \code{RuntimeWarning} per call-site label per process, deduped via the
module-level \code{\_WARNED} set, so a notebook user sees the explanation once
rather than on every call.

\begin{gotcha}
The fork path is \emph{off by default} (\code{parallel=False},
\file{type\_assignment.py:584}). And note the subtle reason the bit-identity
tests do \emph{not} certify notebook safety: pytest is not a ZMQ kernel, so
\code{fork\_unsafe\_in\_notebook} returns \code{False} under pytest and those
tests run on the \emph{fork} path. They prove serial and parallel agree
numerically; they say nothing about whether forking a notebook is safe. Only the
guard protects the notebook. Relatedly, \code{\_ENUM\_WORKER\_STATE} is a module
global populated before the fork and is \emph{not} thread-safe --- do not call
\code{enumerate\_all(parallel=True)} from two threads at once
(\file{type\_assignment.py:657--660}).
\end{gotcha}

% =====================================================================
\section{External tools used in this chapter}\label{sec:ta-tools}

This subsystem leans on three external pieces beyond plain Python. None of NumPy,
SciPy, numba, or networkx appears here. Because the manual assumes you know the
physics but not necessarily the tooling, here is each from the ground up.

\paragraph{SageMath and its symbolic ring \code{SR}.} \term{SageMath} (``Sage'')
is a large open-source mathematics system that bundles many specialised libraries
behind one uniform Python API; the project runs inside a Sage-provided Python
(the \code{MSRJD\_diagrams} conda environment). Its symbolic ring \code{SR} is a
ring of computer-algebra expression trees. \code{SR(x)} coerces a number or
expression into a symbolic expression; methods like \code{.is\_zero()},
\code{.simplify\_full()}, \code{.variables()}, \code{.subs(\{...\})}, and
\code{.simplify\_rational()} operate on it. It is imported in all three core
files (\file{type\_assignment.py:19}, \file{causality.py:16},
\file{symmetry.py:46}). Notably, in \code{type\_assignment.py} the
\emph{semantic} \code{SR(...).is\_zero()} path is \emph{deliberately bypassed} in
favour of the structural string test of \S\ref{sec:ta-zeromask}; in
\code{causality.py}, \code{imag\_part(...)} is Sage's symbolic ``take the
imaginary part'' function.

\paragraph{Sage's graph theory, powered by \code{nauty}.} Sage's \code{DiGraph}
is a directed-graph object; crucially its isomorphism, automorphism, and
canonical-form algorithms are powered by \term{nauty} (``No AUTomorphisms,
Yes?''), Brendan McKay's gold-standard C library for graph canonical labelling
and automorphism-group computation. Given a graph --- optionally vertex-coloured
via a \term{partition} --- nauty computes a canonical form (two graphs are
isomorphic iff their canonical forms match) and the automorphism group (all
relabellings mapping the graph to itself). The partition is how you pass a
colouring: nauty may only map a vertex to another vertex in the \emph{same}
partition cell, which is exactly how the code forbids, say, mapping a noise source
onto an interaction vertex. This subsystem calls into nauty at two points only:
\code{D.automorphism\_group(partition=...)} for the symmetry factor and
compensation (\file{symmetry.py:339}), and
\code{D.canonical\_label(partition=..., certificate=True)} for the dedup signature
(\file{symmetry.py:517}).

\paragraph{SymPy's \code{multiset\_permutations}.} \term{SymPy} is a pure-Python
computer-algebra library, independent of Sage. This subsystem uses exactly one
helper from it, \code{multiset\_permutations} (\file{type\_assignment.py:406}),
the engine of the canonical-leg-matching optimisation of \S\ref{sec:ta-canonical}.

\paragraph{Python standard library.} The rest is stdlib:
\code{itertools.permutations} and \code{itertools.product}
(\file{type\_assignment.py:18}) for orderings and Cartesian products;
\code{collections.Counter} and \code{defaultdict} (\file{symmetry.py:41},
\file{:301}) for tallying multiplicities and grouping edges;
\code{functools.reduce}, \code{operator.mul}, and \code{math.factorial}
(\file{symmetry.py:42--44}) for forming products and factorials;
\code{multiprocessing} (imported lazily, \file{type\_assignment.py:648}) for the
optional fork pool; and \code{IPython.get\_ipython}, \code{sys}, \code{warnings}
(\file{fork\_safety.py}) for the notebook detection and the fallback warning.

% =====================================================================
\section{Gotchas, consolidated}\label{sec:ta-gotchas}

A single reference list of the traps in this subsystem; each is documented in
context above and in the code at the cited lines.

\begin{enumerate}
  \item \textbf{Propagator index order is $(pi, ri)$, not $(ri, pi)$.}
        $G_{\mathrm{ft}}$ is physical-row, response-col, so an edge propagator is
        $G_{\mathrm{ft}}[pi, ri]$. The code \emph{stores}
        \code{propagator\_indices = (ri, pi)} but \emph{looks up} the mask/matrix
        with \code{(pi, ri)}. The historical \code{[ri, pi]} bug
        (\file{type\_assignment.py:490--502}) worked for diagonal couplings and
        silently dropped off-diagonal ones.
  \item \textbf{The zero-check is a string test, not \code{is\_zero()}.}
        \code{str(G\_ft[i,j]) == '0'} (\file{type\_assignment.py:165}) catches
        only literal \code{SR(0)} to avoid multi-second symbolic simplification;
        it can over-generate vanishing diagrams but never produces a wrong number.
  \item \textbf{\code{\_leg\_matchings} emits only the canonical leg ordering.}
        Correct only because the integrator multiplies by $\Sym(\Gamma)$
        afterward; pinned by \code{test\_leg\_matchings\_canonical}.
  \item \textbf{Dedup multiplicity is diagnostic only.} Never multiply it back;
        under Path A the weight is entirely in \code{combinatorial\_factor}
        (\file{symmetry.py:556--568}).
  \item \textbf{\code{diagram\_signature} must stay a complete invariant.} A
        shallow heuristic collided non-isomorphic classes and silently dropped
        integrals; the canonical-form signature cannot
        (\file{symmetry.py:484--498}).
  \item \textbf{\code{external\_wick\_compensation} fails loud on a non-divisor.}
        It raises rather than guess (\file{symmetry.py:391--399}), unlike
        \code{combinatorial\_factor}'s defensive floor-division --- a deliberate
        asymmetry.
  \item \textbf{\code{vertex\_role\_signature} is superseded for weighting.} Its
        depth-1 $\prod N!$ over-divided; the exact compensation replaced it. Kept
        but unused on the production weight path.
  \item \textbf{The fork path is off by default and unsafe in notebooks.}
        \code{parallel=False} (\file{type\_assignment.py:584});
        \code{fork\_unsafe\_in\_notebook} forces serial inside a macOS Jupyter
        kernel. The bit-identity tests run under pytest (not a ZMQ kernel) and
        never certified notebook safety.
  \item \textbf{\code{check\_pole\_structure} is latent.} No \code{pole\_vals} are
        passed in the live pipeline, so only the structural DAG check gates
        diagrams. A pole exactly on the real axis is a hard fail; an inconclusive
        symbolic pole returns \code{passed=True} with conditions.
  \item \textbf{Phase D is not on the live path.} The loop enumerator already
        emits degree-valid prediagrams, so \code{enumerate\_unique\_diagrams}
        skips \code{filter\_prediagrams}; it is exercised by tests and looser
        callers.
\end{enumerate}

% =====================================================================
\section{What you just learned}

This chapter took a bare directed graph and turned it into a set of unique,
causally valid, fully labelled Feynman diagrams, each tagged with its symmetry
factor --- the form the integrators consume. Along the way you saw:

\begin{itemize}
  \item the response/physical edge convention and the load-bearing
        $G_{\mathrm{ft}}[pi, ri]$ index order;
  \item Phase~D's cheap degree pre-screen, and why it is absent from the live
        path;
  \item Phase~E as a backtracking constraint solver, with candidate types,
        external-field permutations, leg matchings, early propagator pruning, and
        the canonical-leg-matching optimisation that keeps the intermediate set
        $15$--$20\times$ smaller with bit-identical output;
  \item the structural string-based zero mask and the performance reasoning
        behind it;
  \item Phase~F's DAG causality check (and its dormant analytic sibling);
  \item Phase~G's symmetry factor via a coloured incidence digraph fed to
        \code{nauty}, the complete-invariant deduplication that closed a real
        silent-wrong-answer bug, and the orbit--stabilizer external Wick
        compensation that hands off to the integrator;
  \item the embarrassingly-parallel fork path and the fork-safety guard that
        makes it survivable on a Mac.
\end{itemize}

The next chapter follows the unique diagrams into the caching layer, which is
what lets a cold enumeration like this one be paid for exactly once.
